\section{Evaluation}

We compare the three architectures along three dimensions: recognition under
session and user variation, deployed system cost, and behavior on continuous
IMU streams.

\subsection{Within-User Recognition Accuracy}

We evaluated session-to-session transfer using five leave-session-out folds.
Each fold used 14 guided sessions for training, two for validation, and four
for testing. Complete sessions remained within one split, and every session
appeared in a test fold exactly once. Training, normalization, early stopping,
and rejection-threshold selection were repeated within each fold using only
its training and validation data.

We report macro-F1 over the four gestures; predictions rejected as null count
as errors. As shown in Table~\ref{tab:within_user}, CNN achieved the highest
mean macro-F1 and the smallest variation across sessions. MLC was 3.3
percentage points lower while executing inside the IMU. Raw-sample HDC was
less accurate and more variable, indicating that its temporal representation
was less robust to session variation.

\begin{table}[t]
    \centering
    \caption{Within-user gesture macro-F1 across five leave-session-out folds.}
    \label{tab:within_user}
    \small
    \begin{tabular}{lcccccc}
        \toprule
        Method & F1 & F2 & F3 & F4 & F5 & Mean $\pm$ SD \\
        \midrule
        MLC & 0.942 & 0.913 & 0.907 & 0.981 & 0.983
            & 0.945 $\pm$ 0.036 \\
        CNN & 0.989 & 0.959 & 0.964 & 0.987 & 0.991
            & 0.978 $\pm$ 0.015 \\
        HDC & 0.699 & 0.557 & 0.719 & 0.750 & 0.732
            & 0.691 $\pm$ 0.077 \\
        \bottomrule
    \end{tabular}
\end{table}

A Python-only diagnostic that encoded the statistical feature vector rather
than raw samples increased HDC macro-F1 to $0.904 \pm 0.039$. This result is
excluded from the deployment comparison, but indicates that input
representation accounts for much of the raw-HDC accuracy gap.

\subsection{Cross-User Transfer}

We evaluate transfer to unseen users using leave-one-participant-out (LOPO)
folds over four participants. In each fold, all data from one
participant are held out, while the remaining participants provide training
and validation data. Training is balanced by participant and gesture class so
that participants with more recordings do not dominate the fitted model. We
report each held-out participant and an equal-participant mean because the
available number of sessions differs across participants.

\begin{table}[t]
    \centering
    \caption{Gesture macro-F1 for transfer to held-out participants. The final
    row reports the equal-participant mean and sample standard deviation.}
    \label{tab:cross_user}
    \small
    \begin{tabular}{lccc}
        \toprule
        Held-out participant & MLC & CNN & HDC \\
        \midrule
        P1 & 0.913 & 0.700 & 0.272 \\
        P2 & 0.787 & 0.696 & 0.327 \\
        P3 & 0.726 & 0.559 & 0.233 \\
        P4 & 0.737 & 0.662 & 0.172 \\
        \midrule
        Mean $\pm$ SD
            & 0.791 $\pm$ 0.085
            & 0.654 $\pm$ 0.066
            & 0.251 $\pm$ 0.065 \\
        \bottomrule
    \end{tabular}
\end{table}

MLC achieved the highest macro-F1 for each held-out participant and an
equal-participant mean of 0.791, compared with 0.654 for CNN and 0.251 for HDC.
In this limited-data setting, the statistical-feature tree transferred more
consistently than either raw-window MCU classifier. All methods declined from
their within-user results. Because participant balancing also caps the number
of training examples per participant, this gap reflects both unseen-user
variation and limited multi-user training data.
Given the small participant sample and unequal recording counts, we treat this
as an exploratory transfer result rather than a population-level estimate.

\subsection{On-Device System Cost}

We compared a sensing baseline with MLC, HDC, and int8 CNN firmware. BLE and
logging were disabled. Current was measured at 4.2\,V over one 60\,s trace
after a 15\,s warm-up; flash and statically allocated RAM were obtained from
the linker reports. Motionless and scripted activity differed by at most
0.02\,mA for all configurations.

\begin{table}[t]
    \centering
    \caption{Power, flash, and static RAM of the firmware configurations.}
    \label{tab:system_cost}
    \small
    \begin{tabular}{
        l
        S[table-format=1.2]
        S[table-format=2.2]
        S[table-format=6.0, group-separator={,}]
        S[table-format=5.0, group-separator={,}]
    }
        \toprule
        {Config.} & {Current (mA)} & {Power (mW)} & {Flash (B)} & {RAM (B)} \\
        \midrule
        Baseline & 1.86 &  7.81 &  63412 & 27288 \\
        MLC      & 1.58 &  6.64 &  63440 & 21184 \\
        HDC      & 3.34 & 14.03 &  76228 & 41208 \\
        CNN      & 4.90 & 20.58 & 170568 & 75568 \\
        \bottomrule
    \end{tabular}
\end{table}

MLC had the lowest measured power and firmware footprint. Its current was
lower than the baseline because classification moved into the IMU and the MCU
no longer drained the raw-data FIFO; this is a deployment-path difference, not
the isolated cost of the tree. HDC occupied an intermediate cost point, while
CNN required the most power, flash, and RAM because of its model, inference
runtime, and tensor arena.

\subsection{Continuous-Stream Behavior}

We replayed guided recordings as continuous streams and measured false
activations on 2.66 hours of held-out free-living data using frozen thresholds.
CNN and HDC advanced 128-sample windows by 64 samples, whereas MLC used its
native 128-sample cadence. With one decision per activation, guided-event
recall was 0.844 for MLC, 0.822 for CNN, and 0.505 for HDC; the corresponding
false-activation rates were 18.78, 84.15, and 728.13 per hour.

Requiring two consecutive decisions reduced false activations to 3.38, 43.35,
and 362.75 per hour, respectively, with guided recall of 0.217, 0.891, and
0.182. These are method-specific operating points rather than an equal-latency
comparison: confirmation adds approximately 1.07\,s for MLC and 0.53\,s for
CNN and HDC. The continuous evaluation therefore exposes a recall--false-
activation trade-off that is not captured by segmented-window macro-F1 alone.
